% Model section of the FISH WiP paper — drafted 2026-09-07 with Claude Fable 5.1 (Anthropic)
\section{The FISH Model}\label{sec:model}

FISH is a hierarchical graph $G=(V,E)$ together with an expansion
operator $\mathcal{X}$. Containment is encoded in the vertices, so $E$
carries only \emph{interaction} edges. Every vertex is a tuple
$v=(t_v,\mathit{id}_v,A_v,Z_v,\mathit{ns}_v)$: its type, a unique
identifier, an attribute tuple, the set of its children in the next
layer, and its ROS~2 namespace. The types induce the ordered hierarchy
$\langle H,EX,N,E,F\rangle$ (host, executor, node, entity, callback) with
layer indices $\ell:\{H,EX,N,E,F\}\to\{0,\dots,4\}$; the sets
$\{Z_v:\ell(t_v)=\ell\}$ partition $V_{\ell+1}$. A node belongs to the
executor vertex of the process that created it, since every trace record
carries the pid that emitted it, and an entity belongs to the callback
group, hence to the executor object, that dispatches it. The attributes are
\[
A_v=\begin{cases}
(\mathit{hostname}) & t_v=H,\\
(\mathit{pid},\ \mathit{mode}\in\{ST,MT\},\ \mathit{threads}) & t_v=EX,\\
(\mathit{label}) & t_v=N,\\
(\mathit{label},\ \mathit{etype},\ \mathit{aspects}) & t_v=E,\\
(\mathit{label},\ \mathit{cb},\ \mathit{phase},\ \mathit{oort}) & t_v=F,
\end{cases}
\]
with $\mathit{etype}\in\{\mathit{tmr},\mathit{sub},\mathit{serv},\mathit{cli},\mathit{wait}\}$.
Here $\mathit{aspects}$ lists the topics and services an entity's
callbacks publish to or consume from, recovered from the link events of
Section~\ref{sec:methodology}; $\mathit{cb}$ is the callback identity;
$\mathit{phase}\in\{\mathit{init},\mathit{data}\}$ separates callbacks
that only ran before the periodic pipeline started from those that ran
after it; and $\mathit{oort}$ flags an \emph{out-of-ROS thread}: a thread
outside the ROS~2 launch tree that exchanges data with ROS~2 nodes over
topics (e.g. a TensorRT worker), observed at the callback layer only.
Namespaces are not a layer: $\mathit{ns}_v$ is the node's ROS~2
namespace, inherited by its entities and callbacks and set-valued for an
executor that serves several namespaces. They only induce a grouping
that presents the user-level abstraction of the architecture: a
depth-one prefix (e.g. \texttt{/perception}) groups the executors
serving it, and a deeper prefix (e.g. \texttt{/perception/tracking})
groups the executors and nodes below it; entities and callbacks follow
their node, and edges between groups are lifted to the group boundary
as between layers.

\emph{a) Edges and expansion:} an interaction edge is a tuple
$e=(u,v,\tau,\mu)$ with endpoints $u,v\in V_4$, a type
$\tau\in\mathcal{T}$ and a payload $\mu$ (topic or service, message
type, and the per-instance measurements of its occurrences). Edges are
established between entities and their callbacks, so $G_3$ and $G_4$
carry the same edges; upward, a vertex inherits every edge between one
of its children and a child of another vertex, whatever its type, so
that with $\pi=p^{\,4-\ell}$
\[
E_\ell=\bigl\{\bigl(\pi(u),\pi(v),\{(\tau,\mu):(u,v,\tau,\mu)\in E_4\}\bigr)
: \pi(u)\neq\pi(v)\bigr\}
\]
and a higher-layer edge is the set of the callback-layer edges it
aggregates. $\mathcal{X}$ is the inverse step: it replaces a vertex by
its children and re-attaches every incident edge to the child endpoints
that carry it,
\begin{align*}
\mathcal{X}_v(G) &= \bigl(V\setminus\{v\}\cup Z_v,\ E'\bigr),\\
E' &= \{(a',b',\tau,\mu):(a,b,\tau,\mu)\in E,\ a'\in\sigma_v(a),\ b'\in\sigma_v(b)\},
\end{align*}
with $\sigma_v(u)=Z_v$ if $u=v$ and $\{u\}$ otherwise. Below the
callback layer, a callback that submits GPU work owns a typed DAG
(Section~\ref{sec:model}c); its incoming edges attach to the entry of
that DAG and its outgoing edges leave from its exit.

\emph{b) Interaction edges:} $\mathcal{T}$ holds the types
$\mathit{msg}$, $\mathit{dep}$, $\mathit{join}$, $\mathit{state}$,
$\mathit{polled}$ and $\mathit{gxf}$, and every edge may carry the flag
$\mathit{external}$.
Let $\mathrm{pub}(f)$ be the topics callback $f$ publishes and
$\mathrm{sub}(T)$ the callbacks of subscriptions to topic $T$. Then
$(f,g,\mathit{msg})\in E_4$ iff $\exists T\in\mathrm{pub}(f):g\in\mathrm{sub}(T)$;
whether the delivery is intra-process or through DDS is an attribute of
$g$, not a different edge. A service call yields a $\mathit{dep}$ pair:
a request edge from the calling callback $f$ to the serving callback $g$
and a response edge from $g$ back to the continuation of $f$ that
handles the reply. Two edge types are \emph{structural inferences} that
recover in-memory hand-offs no ROS layer exposes. If two or more
subscription callbacks of the same node publish the same data topic,
the output is produced by whichever input completes the set: those
callbacks are the completers of a \emph{join} group, which further
contains the node's other subscriptions of the same message type and
the timers that also published the output (the timeout path); members
are tied to the first completer by $\mathit{join}$ edges. The signature
alone does not separate a join from a node that merely forwards whichever
input arrives, so the group is classified from per-instance data: in a
join the completing execution of a callback is an order of magnitude
longer than the deposit-only executions of the same callback (in
Autoware's concatenation node 9.7\,ms against 0.03\,ms) and the output
rate follows the slowest input instead of their sum. Sample-and-hold reads are modelled by $\mathit{state}$ and
$\mathit{polled}$ edges: if a node has subscription callbacks that
consume a data topic but publish none ($S$) and callbacks that publish
data ($P$), every $(s,p)\in S\times P$ with $s$ outside any join group
receives a $\mathit{state}$ edge, inferred, meaning $p$ reads the latest
value deposited by $s$; when that read is observed as a message take
inside $p$'s execution, the edge is a $\mathit{polled}$ edge instead.
Both carry the measured data age rather than a hop latency. $\mathit{gxf}$ edges are the GXF-internal hops of NITROS
graphs, every ingress of a codelet group feeding every egress. Finally,
an untraced endpoint (a topic with traced subscribers but no traced
publisher, or vice versa) and every out-of-ROS thread are represented
by an external vertex; the edges into and out of it are marked
$\mathit{external}$, and when lifted the vertex becomes an external
entity of that layer, which is how a composed system meets its
environment.

\emph{c) GPU sub-model:} a callback $f$ that submits GPU work owns a set
of typed DAGs, one per distinct \emph{signature} observed over all its
invocations, where an invocation is one execution window of $f$ that
contains CUDA runtime calls. A typed DAG $D_f^k=(N,A)$ has host calls
($\mathit{cpu}$), kernels ($\mathit{sm}$) and copies ($\mathit{ce}$) as
nodes, and $\mathit{launch}$ (a host call to the work it produced, by
correlation id), $\mathit{cpu\_order}$ (consecutive host calls) and
$\mathit{stream\_fifo}$ (consecutive work on one stream) as edges; its
signature is its node sequence. Over this set FISH classifies shapes
and phases (the dominant steady shape against initialisation, warm-up
and transient variants) and derives the callback's \emph{conditional
graph}: the common backbone of all signatures with the observed
branches and their probabilities, which is the input-dependent
structure a GPU-aware analysis has to bound.

\emph{d) Correctness:} correctness is stated, and checked, on the graph
as extracted, before any grouping or filtering is applied. Let $\chi$ map every vertex-creating ROS~2 API call to
the trace records it must emit and the vertices it must create: a node
constructor yields $1N+7E+7F$, because it also creates the six parameter
services and the time-source subscription, whereas
\texttt{create\_service} yields $1E+1F$ and \texttt{create\_publisher}
yields no vertex at all, only a publish aspect on its caller. $\chi$ is
tabulated once for the ROS~2 API and for the wrapper libraries that call
it. For a run, FISH locates the sources of every binary its traced
processes mapped, scans them for API calls and sums the expected records
and vertices; correctness is then the two-phase containment
\[
\chi_r(A)\subseteq\mathcal{R}\ \ \text{(C1)},\qquad
\chi_v(A)\subseteq V\ \ \text{(C2)},
\]
where $A$ is the set of executed API calls, $\mathcal{R}$ the trace and
$V$ the FISH graph: C1 is instrumentation completeness, every expected
record is in the trace, and C2 is extraction completeness, every
expected vertex is in the graph, checked node by node. The views used for presentation are lossy by design: each
applies a declared filter predicate from a catalogue, so what a view
omits is always attributable to a filter, never to the extraction. The single exclusion before extraction is by
definition, not by filtering: a process that FISH killed and relaunched
is an artefact of tracing rather than a part of the application, and its
records are removed from $\mathcal{R}$.
